% ===================================================================
%  TABULKA č. 2 — Lidské tělo. — Přestávka. — Hry.
%  Traduction 2026 d'apres texto/io, controlee sur texto/fr, texto/ru
%  et texto/uk. Consigne : voir texto/cs/10-tabulka-01.tex.
%
%  LE TABLEAU DU CORPS EST LE PREMIER BANC D'ESSAI DE CETTE COLONNE,
%  et il montre tout de suite ce qui la separe de l'ukrainien, sa
%  seule voisine slave dans le volume. Les deux langues sont slaves,
%  mais le tcheque a garde des mots que l'ukrainien a remplaces et
%  inversement : « čelo » le front et « oko » l'oeil sont communs,
%  mais « tvář » la joue, « brada » le menton, « rty » les levres,
%  « koleno » le genou, « kotník » la cheville, « lýtko » le mollet
%  ne se lisent pas de la meme facon a Kiev.
%
%  ET LES JEUX N'ONT PAS DE NOM INTERNATIONAL, comme l'ido le dit
%  lui-meme dans la note du feuillet 17. Le tcheque a les siens :
%  « slepá bába » pour le colin-maillard, « na schovávanou » pour
%  cache-cache, « káča » la toupie, « kuličky » les billes,
%  « švihadlo » la corde a sauter, « obruč » le cerceau.
% ===================================================================

%%K t02-apar-1 apar
\VUsaut{4.0mm}
\VUtitre{15.0pt}{60}{TABULKA č. 2}
\VUsaut{2.0mm}
\VUtitre{11.4pt}{0}{Lidské tělo. --- Přestávka.}
\VUtitre{11.4pt}{0}{Hry.}

%%K t02-tit-0 sub
\VUtitre{13.2pt}{0}{Lidské tělo.}
\VUsaut{2.4mm}
\VUcentre{10.2pt}{0}{\textit{(Jednoduchá lekce přírodopisu.)}}
\VUsaut{3.0mm}

%%K t02-01-1 p
1. --- Hlavní části lidského těla jsou: \VUgras{hlava,} \VUgras{trup} a
\VUgras{údy.}

%%K t02-tit-1 sub
\VUpk{10.2pt}{40}{Hlava.}
\VUsaut{3.0mm}

%%K t02-02-1 p
2. --- Hlava zahrnuje dvě části: \VUgras{lebku,} která obsahuje \VUgras{mozek} a kterou
pokrývají \VUgras{vlasy,} a \VUgras{obličej.}

%%K t02-03-1 p
3. --- Na našem obrázku nevidíme ani lebku, ani mozek; ale vidíme velmi zřetelně důležité
části obličeje.

%%K t02-04-1 p
4. --- Zde je nejprve \VUgras{čelo,} které prozrazuje mocný rozum, myšlení stále činné.
\VUgras{Spánky} jsou zcela volné. \VUgras{Oko} je živé a široce otevřené. Husté \VUgras{obočí}
a dlouhé \VUgras{řasy} je chrání.

%%K t02-05-1 p
5. --- Zde je dále \VUgras{nos} a \VUgras{nozdry.} Nos nám slouží k čichání a k dýchání. Po
každé straně nosu jsou \VUgras{tváře,} dále \VUgras{uši.} Spodní část obličeje tvoří
\VUgras{ústa} a \VUgras{brada.}

%%K t02-06-1 p
6. --- Nevidíme je příliš dobře, protože \VUgras{knír} a \VUgras{licousy} nám je zakrývají.
Ale víme, že ústa mají dva \VUgras{rty,} \VUgras{horní čelist} a
\VUgras{dolní čelist} se svými \VUgras{zuby,} a \VUgras{jazyk.} Ústy mluvíme a jíme.
Jazykem chutnáme pokrmy.

%%K t02-07-1 p
7. --- Oko, ucho, nos a ústa jsou, stejně jako prsty, orgány pěti smyslů:
\VUgras{zraku,} \VUgras{sluchu,} \VUgras{čichu,} \VUgras{chuti,} \VUgras{hmatu.}

%%K t02-08-1 p
8. --- Hlava je tedy jednou z nejzajímavějších částí lidského těla.

%%K t02-tit-2 sub
\VUsaut{4.4mm}
\VUpk{10.2pt}{40}{Trup.}
\VUsaut{3.0mm}

%%K t02-09-1 p
9. --- \VUgras{Krk} spojuje hlavu s trupem. Zahrnuje \VUgras{hrdlo} a \VUgras{šíji.}

%%K t02-10-1 p
10. --- Trup zahrnuje také mnoho podstatných orgánů. V \VUgras{hrudi} jsou
\VUgras{plíce,} kterými dýcháme, a \VUgras{srdce,} které je středem života. Když
srdce přestane tepat, živý tvor umírá.

%%K t02-11-1 p
11. --- \VUgras{Žebra} drží hruď.

%%K t02-12-1 p
12. --- Ostatní části trupu jsou \VUgras{bok,} \VUgras{záda,} \VUgras{ramena} a
\VUgras{břicho.} Léháme si na bok nebo na záda, abychom spali, a náklady nosíme na
svých ramenou.

%%K t02-13-1 p
13. --- V břiše jsou \VUgras{žaludek} a \VUgras{střeva.} Slouží nám k trávení pokrmů.

%%K t02-tit-3 sub
\VUsaut{4.4mm}
\VUpk{10.2pt}{40}{Údy.}
\VUsaut{3.0mm}

%%K t02-14-1 p
14. --- Máme čtyři \VUgras{údy}: dvě \VUgras{paže} a dvě \VUgras{nohy.} Pažemi bereme, držíme,
zvedáme a sbíráme předměty; nohama stojíme, chodíme a běháme.

%%K t02-15-1 p
15. --- \VUgras{Rameno} spojuje paži s tělem. Velmi mohutný \VUgras{sval,} biceps, pohybuje
paží, kterou \VUgras{loket} spojuje s \VUgras{předloktím.} \VUgras{Zápěstí} tvoří kloub
\VUgras{ruky,} která končí \VUgras{prsty} a \VUgras{nehty.} Na ruce rozeznáváme \VUgras{žílu}
(a tepnu), v níž proudí \VUgras{krev.}

%%K t02-16-1 p
16. --- Části nohy jsou: \VUgras{stehno,} \VUgras{koleno} a \VUgras{podkolení,}
\VUgras{lýtko,} jehož \VUgras{maso} je pokryto \VUgras{kůží,} \VUgras{kotník} a
\VUgras{chodidlo.} \VUgras{Kosti} nohy jsou velmi silné a velmi pevné. Na konci
chodidla jsou \VUgras{prsty u nohy.} Ostatní části jsou \VUgras{ploska} a
\VUgras{pata.}

%%K t02-17-1 p
17. --- Takový je téměř úplný popis lidského těla.

%%K t02-17-2 p
\textit{Poznámka.} --- \textit{Pomocí obratu: « bolí mě... (hlava atd.) » bude moci učitel
použít celou tuto slovní zásobu znovu z hlediska dobrého nebo špatného}
\VUgras{tělesného stavu.}

%%K t02-tit-4 sub
\VUsaut{5.0mm}
\VUpk{10.2pt}{40}{Tělocvik.}
\VUsaut{2.4mm}
\VUcentre{10.2pt}{0}{\textit{(Vypravování jednoho ze žáků.)}}
\VUsaut{3.0mm}

%%K t02-18-1 p
18. --- Abychom byli ohební, mrštní a zdraví, musíme cvičit své nohy a své paže. Proto si
hrajeme a pěstujeme \VUgras{tělocvik.}

%%K t02-19-1 p
19. --- Během týdne se tato dvě cvičení konají na dvoře lycea (viz tabulku č. 1). Ale ve
čtvrtek a v neděli chodíme na velký trávník, kde máme místo k běhání a k zábavě.

%%K t02-20-1 p
20. --- Naši učitelé a naši rodiče nás tam někdy doprovázejí; ale cvičení řídí opravdu
\VUgras{učitel tělocviku} \textsuperscript{(12)}.

%%K t02-21-1 p
21. --- Na trávníku, vlevo od vchodu, je \VUgras{tělocvičné nářadí} \textsuperscript{(1)};
šplháme po \VUgras{žebříku} \textsuperscript{(2)}, děláme převraty na \VUgras{kruzích}
\textsuperscript{(3)} a na \VUgras{hrazdě} \textsuperscript{(4)}, vytahujeme se rukama až na
vrchol \VUgras{provazového žebříku} \textsuperscript{(5)} nebo \VUgras{uzlovaného lana}
\textsuperscript{(6)}.

%%K t02-22-1 p
22. --- Zatím naši kamarádi skáčou přes \VUgras{dřevěného koně} \textsuperscript{(7)} nebo
přes \VUgras{bradla} \textsuperscript{(11)}. Jiní
\VUgras{boxují} \textsuperscript{(8)} nebo \VUgras{šermují} \textsuperscript{(9)} s
maskou a \VUgras{fleretem.} Konečně jiní \VUgras{zvedají činky} \textsuperscript{(10)}.

%%K t02-tit-5 sub
\VUsaut{4.6mm}
\VUpk{10.2pt}{40}{Hry.}
\VUsaut{3.0mm}

%%K t02-23-1 p
23. --- Kolem cvičenců si chlapci a děvčata hrají různé \VUgras{hry.} Děvčata se baví se svou
\VUgras{obručí} \textsuperscript{(14)}; skáčou přes \VUgras{švihadlo} \textsuperscript{(13)}
nebo zvou své bratry a bratrance na partii \VUgras{kroketu} \textsuperscript{(15)}. Mají
velkou radost, když svou \VUgras{dřevěnou paličkou} odpálí daleko \VUgras{kouli} protivníka
právě ve chvíli, kdy si myslí, že snadno projde brankou!

%%K t02-24-1 p
24. --- Chlapci dávají přednost hrám více svalovým. Rádi hrají \VUgras{kuželky}
\textsuperscript{(16)}, které obratně porážejí \VUgras{koulí} \textsuperscript{(17)}.
Nepohrdají ani hrou s \VUgras{káčou} \textsuperscript{(18)} a s \VUgras{kuličkou a kalíškem}
\textsuperscript{(19)} nebo dokonce \VUgras{žabí hrou} \textsuperscript{(20)}. Ale jejich
nejoblíbenější hry jsou opravdu \VUgras{kuličky} \textsuperscript{(21)} a \VUgras{míč o zeď}
\textsuperscript{(22)}, protože zeď
\VUgras{skleníku} \textsuperscript{(26)} je velmi vhodná k tomu, aby lehký míček
odrážela.

%%K t02-25-1 p
25. --- Malý Jan, který chodí na svých \VUgras{chůdách} \textsuperscript{(24)}, je velmi
obratný a umí velmi dobře udržet rovnováhu. Blízko něho si někteří kamarádi hrají
\VUgras{na schovávanou} \textsuperscript{(25)} v tom skleníku a za
\VUgras{živým plotem.} Jiní dávají přednost \VUgras{horké ruce}
\textsuperscript{(28)} nebo \VUgras{badmintonovému míčku} \textsuperscript{(29)}, který létá z
jedné \VUgras{rakety} na druhou.

%%K t02-noto-1 noto
\VUnotes{143.2mm}{%
(*) Dětské nebo chlapecké hry mají často jména čistě národní. Snažili jsme se pojmenovat je v
Idu co nejlépe, buď podle samotného úkonu, jímž se vyznačují, buď volbou národního jména v té
či oné zemi, není-li příliš nepatřičné. Např.:
\VUgras{hra na sud} (německy a francouzsky) je \VUgras{žabí hra}
\textsuperscript{(20)} (anglicky), v níž se hráči snaží vhodit své kulaté kotoučky ústy kovové
žáby, která stojí s otevřenou tlamou na proděravěné bedně (na sudu).
\VUgras{Skok přes rameno} \textsuperscript{(30)} se od \VUgras{skoku přes záda} liší
jen tímto: aby přeskočili přes hlavu, opírají si hráči ruce o ramena svého druha, kdežto při «
\VUgras{skoku přes záda} » si opírají ruce jen o jeho záda. --- Anebo jsou někteří z účastníků
skloněni, zatímco ostatní budou přeskakovat přes jejich záda opírajíce si ruce o rameno; první
musí snést náraz, aniž se podlomí, druzí musí skočit, aniž ztratí rovnováhu (Viz samotnou
tabulku).%
}

%%K t02-26-1 p
26. --- Uprostřed trávníku jsou dva stromy. U paty jednoho z nich uspořádalo sedm nebo osm
chlapců partii \VUgras{skoku přes rameno} \textsuperscript{(30)} (*). Ne ve všech zemích se
tato hra zná; ale všichni vědí, co je \VUgras{skok přes záda,} který chlapci tentokrát
nehrají.

%%K t02-tit-6 sub
\VUsaut{5.0mm}
\VUpk{10.2pt}{40}{Hry na volném vzduchu.}
\VUsaut{3.4mm}

%%K t02-27-1 p
27. \VUgras{Slepá bába} \textsuperscript{(31)} je také velmi zábavná; děvčata a chlapci si
střídavě dávají na oči \VUgras{šátek} a chytají neobratné, kteří se nechají polapit.

%%K t02-28-1 p
28. --- Uprostřed trávníku, zde je hra velmi francouzská, ale poněkud prudká: je to
\VUgras{medvěd} \textsuperscript{(32)}, mnohem hlučnější než tak klidná hra s
\VUgras{drakem} \textsuperscript{(33)} nebo dokonce než anglická hra \VUgras{tenis}
\textsuperscript{(34)}.

%%K t02-29-1 p
29. --- Tento poslední sport je radostí starších žáků. Sami naši učitelé jej rádi pěstují.
Vyžaduje velkou obratnost spolu s mrštností a velmi se mi líbí. Přenechávám děvčatům hry na
\VUgras{houpačce} \textsuperscript{(35)}.

%%K t02-30-1 p
30. --- Nemám rád jinou anglickou hru, která by se snad mohla stát nebezpečnou.

%%K t02-30-1b p
Chci mluvit o \VUgras{kopané} \textsuperscript{(36)}, která těší mého staršího bratra. Dávám
přednost naší staré \VUgras{hře na zajatce} \textsuperscript{(38)}, stejně zdravé a opravdu
méně surové.

%%K t02-tit-7 sub
\VUsaut{4.6mm}
\VUpk{10.2pt}{40}{Cyklistika a balon.}
\VUsaut{3.0mm}

%%K t02-31-1 p
31. --- Také rád objíždím trávník na \VUgras{jízdním kole} \textsuperscript{(37)}.

%%K t02-32-1 p
32. --- Příští rok se naučím \VUgras{jezdit na koni} \textsuperscript{(39)}. Jak je krásný
kůň, který půvabně kráčí nebo klusá a který cvalem skáče přes překážky!

%%K t02-33-1 p
33. --- V létě se učíme \VUgras{plavat} \textsuperscript{(40)}. S rozkoší se noříme a koupeme
v čisté a svěží vodě řeky. Někdy nakonec vypouštíme papírový \VUgras{balon}
\textsuperscript{(41)}, který velebně stoupá do ovzduší, jakmile je naplněn teplým vzduchem.

%%K t02-34-1 p
34. --- Je také mnoho jiných her, ale nedokončil bych jejich výčet, a z tohoto shrnutí se
chápe, že se ve čtvrtek na našem krásném trávníku příliš nenudíme.

%%K t02-34-2 p
N.-B. --- \textit{Učitel snadno doplní tuto kapitolu, když podrobněji popíše každou z
nejdůležitějších her.}
